%% alg-resolve.tex -- Algorithm: the client-side read path.
%%
%% Short on purpose. This is the only CALIBER library code on the serving path,
%% so it is also the only place where the control plane can cost a production
%% request anything -- which is what E1 measures. Showing it makes the size of that
%% cost surface auditable rather than asserted.

\begin{algorithm}[tbp]
\caption{The client-side read path --- the only \sys code on the serving path, and
therefore the entire surface E1 measures. Late binding lives on line~2: the client
names an alias, never a version, which is what decouples remediation from
deployment. A TTL cache can satisfy line~2 without an HTTP call. Line~5 is the property that makes the resolver's failure mode
acceptable: a resolver that cannot reach the control plane serves the last known
version rather than failing the call.}
\label{alg:resolve}
\psig{resolve}
\pin{$a$}{an asset identity; $\alpha$, an alias such as \code{@prod}}
\pout{the definition to execute now}{}
\begin{pseudo}
\pline{$k \gets (a, \alpha)$}
\pline{$v \gets$ \kw{lookup} \pid{live\_target}($k$)
       \cmt{late binding; TTL cache permitted}}
\pline{\kw{if} $v \neq \bot$ \pthen}
\pline{\I cache($k$) $\gets v$; \kw{return} definition($v$)}
\pline{\kw{else if} cache($k$) $\neq \bot$ \pthen}
\pline{\I \kw{return} definition(cache($k$))
        \cmt{degrade, do not fail}}
\pline{\kw{else}}
\pline{\I \kw{raise} unresolved($k$)
        \cmt{explicit, never a silent default}}
\end{pseudo}

\vspace{0.5ex}
\pinv{a client never names a version. Consequently a release changes behaviour
  with no client deployment, and a rollback is observed by the next call rather
  than by the next deploy.}
\end{algorithm}
